\documentclass[10pt]{article}
\usepackage[letterpaper,margin=.45in]{geometry}
\usepackage{amsmath,amssymb,multicol}
\pdfinfoomitdate=1
\pdftrailerid{}
\pagestyle{empty}
\setlength{\parindent}{0pt}
\setlength{\parskip}{2pt}
\setlength{\jot}{0pt}
\let\bigsum\sum
\renewcommand{\sum}{\bigsum\nolimits}
\setlength{\columnsep}{14pt}
\setlength{\abovedisplayskip}{3pt}
\setlength{\belowdisplayskip}{3pt}
\newcommand{\fl}[1]{\left\lfloor #1\right\rfloor}
\newcommand{\rng}[2]{\{#1,\ldots,#2\}}
\begin{document}
\begin{center}
{\large\bfseries A completely specified Diophantine consistency equation}\\[-1pt]
{\small Fifteen strictly positive integer unknowns; all fixed constants defined on this page.}
\end{center}
\vspace{-7pt}
\begin{multicols}{2}
\small
The unknowns are
\[
 J,F,\alpha,z,f,h,i,T,s,w,\tau,\eta,\zeta,y,\gamma\in\mathbb Z_{>0}.
\]
Use the following polynomial abbreviations (no extra unknowns):
\[
\begin{gathered}
q=\beta J+1,\quad C=q-F-\alpha,\quad X=wq,\\
Y=sq^3,\quad E=XY,\quad k=\eta+\zeta,\quad c=kY+\eta,\\
a=E+Y,\quad H=4a+3,\quad\Delta=a^2+H,\\
D=X+ac+\gamma H,\\
R=((q-1)(q-F)+\alpha)(q^2-1)+(M_C+qM_F)J,\\
V=c(Tf-1)-Rf^2.
\end{gathered}
\]
The equation is
\[
\boxed{N_0N_1N_aN_kN_tN_s-1=0,}
\]
where
\begin{align*}
N_0&=\tau^2-(XY^2)(XY^2+1)k^2,\\
N_1&=D^2-\Delta c^2,\\
N_a&=\Delta(f^2-1)(V^2-y^2)+y^2,\\
N_k&=k-hE-R,\\
N_t&=(K+w)C+(q-F)-z(q-1),\\
N_s&=f^2-\Delta(ic^2)^2.
\end{align*}

\textbf{Fixed coefficients.}
All sums below have the displayed fixed finite limits. Division and
floor are used only to define constants; the equation above is polynomial.
Put
\[
\begin{gathered}
\rho=2^{89},\quad B=\rho^{4831},\quad\beta=B-1,\\
\lambda=2^{29},\quad\theta=2^{42},\quad Q=2^{40},
\end{gathered}
\]
\[
\begin{gathered}
G_m(x)=\sum_{j=0}^{m-1}x^j,\\
\mu=510+2^{17}+2^{26}+2^{28}+\lambda(2^{44}-1).
\end{gathered}
\]
Let
\begin{align*}
\mathcal S={}&\rng0{456}\cup\rng{495}{504}\cup\{507,508\}\\
 &\cup\rng{517}{537}\cup\rng{539}{549}\\
 &\cup\rng{567}{583}\cup\{1203,1807\}.
\end{align*}
Then
\begin{align*}
M_C={}&B-1-\sum_{e\in\mathcal S\setminus\{1\}}\rho^e-2\rho^{507},\\
M_F={}&B-1+4+\mu(\rho^{2409}+\rho^{3015})\\
 &+\rho^{528}G_{10}(\rho)+\rho^{585}G_{11}(\rho)+\rho^{1203}+\rho^{1807}\\
 &+(2^{87}-1)\rho^{4823}G_4(\rho)+(2^{38}-1)\rho^{4827},\\
K={}&\rho^{33}+\rho^{68}+\rho^{664}+\rho^{1268}+\rho^{4877}\\
 &+(\rho^{2409}+\rho^{3015})\Phi(1/\rho),
\end{align*}
where the rational intermediate expression is
\begin{align*}
\Phi(x)={}&\sum_{i=0}^{456}c_i x^{r_i}\\
 &+\lambda\bigl[(x^{539}+2^{21}x^{517})G_{11}(2x)\\
 &\qquad +(2^{11}x^{528}+2^{32}x^{495})G_{10}(2x)\\
 &\qquad +\theta(x^{539}+x^{1203}+x^{1807})\bigr]\\
 &+\bigl[(2^9+\lambda(2+2^{33}))x^{567}\\
 &\qquad +(2^{18}+\lambda(2^{22}+2^{33}))x^{575}\bigr]G_8(2x)\\
 &+(2^{27}+\lambda(2^{11}+2^{32}))x^{583}.
\end{align*}
The resulting $\beta,K,M_C,M_F$ are fixed positive integers.

\textbf{The constants $r_i,c_i$.}
First define the following 457 fixed quadruples
$t_i=(w_i,n_i,e_i,s_i)$.
From the integer $\mathcal A$ printed below, for $0\le j<430$ put
\[
\begin{gathered}
b_j=\fl{\mathcal A/1024^j}\bmod1024,\quad u_j=b_j\bmod256,\\
v_j=\fl{b_j/256}\bmod2,\quad d_j=\fl{b_j/512}.
\end{gathered}
\]
Here $a\bmod b=a-b\fl{a/b}$. Set
\[
t_{2+j}=\bigl((1-d_j)(512+u_j),\ 512+j,\ d_j(256+u_j),\ v_j\bigr).
\]
The remaining quadruples are
\[
\begin{array}{c|c|l}
\text{index}&\text{range}&t_i\\\hline
0&& (2,3,3,1)\\
1&& (5,0,0,0)\\
432+j&0\le j<2&(0,942+j,0,942+j)\\
434+j&0\le j<2&(0,j,0,j)\\
436+j&0\le j<4&(0,(j\bmod2)+942\fl{j/2},0,6)\\
440&&(2,2,2,0)\\
441&&(3,4,4,698)\\
442&&(4,4,4,0)\\
443+j&0\le j<3&(1,6,1,2+j)\\
446+3j+k&0\le j,k<3&(j^2,5,k,5)\\
455&&(256,0,0,512)\\
456&&(0,0,512,512)
\end{array}
\]
Define, for these fixed indices,
\begin{align*}
\kappa_i&=2^{30}w_i+2^{20}n_i+2^{10}e_i+s_i,\\
\ell_i&=2w_i+2^{11}n_i+2^{22}e_i+2^{32}s_i,\\
r_0&=0,\quad r_1=1,\\
r_i&=2+\sum_{j=2}^{456}\fl{\frac{Q+\kappa_i-\kappa_j-1}{Q}}
 \quad(2\le i\le456),\\
\delta_{ij}&=\fl{\frac{Q-(i-j)^2}{Q}},\\
O&=1+255\,2^9+255\,2^{18}+2^{27}+\lambda,\\
c_i&=O+\lambda\ell_i-255\,2^9\delta_{i,455}-255\,2^{18}\delta_{i,456}\\
 &\qquad -2^{27}(\delta_{i,455}+\delta_{i,456}).
\end{align*}
All subscripts in these definitions are bound constant indices, not
additional Diophantine unknowns.

\textbf{The integer $\mathcal A$.} Concatenate the following lines in order:
\begingroup
\fontsize{8}{9.1}\selectfont\ttfamily
%% Generated by tools/build_one_page_equation.py; concatenate these lines.
56118335615737617036884126988661458925244389493947058121715971\par
01282916856878425443037240736299803427824619599002153617345708\par
28900551636239650126542082909466175026205487572409217714429639\par
74793862932436776448290221360877005574415446786244745611959287\par
59373112491605173997640015648449178599058708778363556006068759\par
35678639237640904034915105614223460060014840095836650944470526\par
30060761226207037490030688247103632272616121659441631458841739\par
88053038481633961651885200391264203970434788452535325946647905\par
26112698210284600353636434795855158687861272006223054621338052\par
08060289634159601137749386030632145140305682244017830050480315\par
09861375430770373735746471163208368830937059811926841173596438\par
99567325570145003948256119536072783258394320179887629332814913\par
83742105792556854838396752391719155979103190193016087490398924\par
39017272426043275053833387572790002456817641447316642739994186\par
97003064581790026300137955120486268219145788062206957924754981\par
34155288871384611877132683517463256922006640654344797535262481\par
90435827191666264404447576661943567063783141522216477014676593\par
40310593002091093201099250164766043944383508572504290831392543\par
80006501642309702444745828581986468804777321875277055695293159\par
15172377582008682228504985026248220701674081629382752859968808\par
927940829152477998549834419730978516012538226242487133\par

\endgroup
\medskip
\normalfont\footnotesize
Absence of positive integer solutions is equivalent to $\mathrm{Con}(\mathrm{ZFC})$
by the inherited native/Pell argument and the checked P69--P67 reduction.
The full native calibration and its internal ZFC formalization remain open
Lean obligations. This page specifies the equation, not its proof.
\end{multicols}
\end{document}
